\medskip

\begin{enumerate}
	\item 
	
$\bullet~~$On sait que $\displaystyle\lim_{x \to 0} x \ln x = 0$, donc $\displaystyle\lim_{x \to 0} f(x) = 1$.

$\bullet~~$Comme $\displaystyle\lim_{x \to + \infty} x = + \infty$ et $\displaystyle\lim_{x \to + \infty} \ln x = + \infty$, on a par produit de limites $\displaystyle\lim_{x \to + \infty} x \ln x = + \infty$ et enfin $\displaystyle\lim_{x \to + \infty} f(x) = + \infty$.
	\item \begin{enumerate}
		\item La fonction $f$ est continue et dérivable sur $]0~;~+\infty[$. En utilisant le formule de la dérivée d'un produit, $\forall x \in ]0~;~+\infty[$, $f'(x)=1\times \ln(x)+x\times \dfrac{1}{x}=\ln(x)+1$
		\item $f'(x)\geqslant 0 \iff \ln(x)+1 \geqslant 0 \iff \ln(x) \geqslant -1 \iff x\geqslant \e^{-1}$.
		
		Dans le tableau : $f\left(\e^{-1}\right)= \e^{-1} \ln(\e^{-1})+1=-\e^{-1}+1=1-\e^{-1}>0$
\begin{center}
\begin{pspicture}(8,3)
\psframe(8,3)\psline(0,2)(8,2)\psline(0,2.5)(8,2.5)
\psline(2,0)(2,2.5)\psline(2.05,0)(2.05,2.5)
\uput[u](1,2.4){$x$}\uput[u](2.025,2.4){$0$}\uput[u](5,2.4){$\text{e}^{-1}$}\uput[u](7.5,2.4){$+ \infty$}\uput[u](6.5,2.4){$1$}
\rput(1,2.25){$f'(x)$}\uput[u](3.5,1.9){$-$}\uput[u](5,1.9){$0$}\uput[u](6.5,1.9){$+$}
\rput(1,1){$f(x)$}\uput[d](2.2,2){1}\uput[u](5,0){$1 - \text{e}^{-1} \approx 0,632$}
\uput[d](7.5,2){$+ \infty$}\rput(6.5,1){1}
\psline{->}(2.2,1.5)(4.5,0.5)\psline{->}(5.5,0.5)(7.5,1.5)
\end{pspicture}
\end{center}
		\item $f(1)=1$. D'après les variations de la fonction $f$,
\begin{itemize}
	\item $\forall x \in ]0~;~\e^{-1}]$, $f(x)\in [1-\e^{-1}~;~1[$ ;
	\item $\forall x \in [\e^{-1}~;~1]$, $f(x) \in [1-\e^{-1}~;~1]$.
\end{itemize}

Donc $\forall x \in ]0~;~1]$, $f(x)\in \left[1-\e^{-1}~;~1\right]$. Or $1 - \e^{-1} > 0$ donc $f(x)\in ]~0~;~1]$ pour tout réel $x$ dans $]0~;~1]$.
		\end{enumerate}
	\item 
	\begin{enumerate}
		\item L'équation de la tangente $(T)$ à $\mathcal{C}_f$ au point d'abscisse 1 a pour équation :

$y=f'(1)(x-1)+f(1)$

		 Avec $f'(1)=1$ et $f(1)=1$, on en déduit l'équation de $(T)~:~y=x-1+1=x$, soit 

$M(x~;~y) \in (T) \iff y = x$.
		\item Nous savons que $\forall x \in ]0~;~+\infty[$, $f'(x)=\ln(x)+1$.

		La fonction $f'$  est continue et dérivable sur $]0~;~+\infty[$, et $f''(x)=\dfrac{1}{x}>0$, donc $f$ est convexe sur $]0~;~+\infty[$.
		\item La courbe représentative d'une fonction convexe est au-dessus de toutes ses tangentes. En particulier $\mathcal{C}_f$ est au-dessus de $(T)$. Donc $\forall x \in ]0~;~+\infty[$, $f(x)\geqslant x$.
	\end{enumerate}
	\item Soit $\left(u_n\right)$ la suite définie par $u_0 \in ]0~;~1]$ et $u_{n+1}=f(u_n)$.
	\begin{enumerate}
		\item Montrons par récurrence que $0 < u_n <1$ pour tout entier naturel $n$.
		
\emph{Initialisation} : on a $u_0 \in ]0~;~1[$, soit $0 < u_0 < 1$ : l'encadrement est vrai au rang $0$ ;

\emph{Hérédité} : soit $n \in \N$ et supposons que $0 < u_n < 1$.

On a vu à la question \textbf{2. b.} que si $0 < x < 1$, alors $0 < f(x) < 1$.

Donc si $0 < u_n < 1$, alors $0 < f\left(u_n\right) < 1$, soit $0 < u_{n+1} < 1$.

\emph{Conclusion} : l'encadrement est vrai au rang 0 et s'il est vrai au rang $n \in \N$, il l'est aussi au rang $n + 1$ : d'après le principe de récurrence pour tout entier naturel $n$, on a : $0 < u_n < 1$.
		

		\item D'après la question \textbf{3.c}, pour tout réel $x$ positif, $f(x)>x$.

De plus pour tout entier naturel $n$, $u_n>0$

Donc, $\forall n \in \N$, $f(u_n)>u_n$ donc $u_{n+1}>u_n$. La suite $\left(u_n\right)$ est donc croissante.
		 \item La suite $\left(u_n\right)$ est croissante et majorée par 1. Donc d'après le théorème de convergence monotone, la suite $\left(u_n\right)$ converge vers une limite finie notée $l$.
	\end{enumerate}
\end{enumerate}

\pagebreak
